\chapter{ความหลากหลายและโครงสร้างของจุลินทรีย์ในดิน}
\label{chap:ch01}

% ==========================================================
% แผนบริหารการสอนประจำบทที่ 1
% ==========================================================
\section*{แผนบริหารการสอนประจำบทที่ 1}
\addcontentsline{toc}{section}{แผนบริหารการสอนประจำบทที่ 1}

\noindent\textbf{หัวข้อเนื้อหาประจำบท}
\begin{enumerate}
    \item ดินในฐานะระบบนิเวศและแหล่งอาศัยของสิ่งมีชีวิต
    \item การจำแนกและลักษณะของกลุ่มจุลินทรีย์หลักในดิน (แบคทีเรีย รา แอกทิโนมัยซีท อาร์เคีย และโปรโตซัว)
    \item สถาปัตยกรรมเม็ดดิน (Soil Aggregates) และแหล่งอาศัยระดับไมโคร
    \item ปัจจัยทางกายภาพ เคมี และชีวภาพที่ควบคุมความหลากหลายของจุลินทรีย์ดิน
    \item การประเมินชีวมวลและกิจกรรมของจุลินทรีย์ในดิน
\end{enumerate}

\noindent\textbf{วัตถุประสงค์เชิงพฤติกรรม}
\begin{enumerate}
    \item อธิบายลักษณะโครงสร้าง บทบาท และเกณฑ์การจำแนกกลุ่มจุลินทรีย์หลักในดินเกษตรได้อย่างถูกต้องและครบถ้วน
    \item วิเคราะห์ความสัมพันธ์ระหว่างโครงสร้างเม็ดดินกับแหล่งอาศัยระดับจุลภาคของประชากรจุลินทรีย์ดินได้อย่างมีเหตุผล
    \item อธิบายอิทธิพลของค่า pH ความชื้น และปริมาณอินทรียวัตถุในดินต่อการเปลี่ยนแปลงโครงสร้างชุมชนจุลินทรีย์ได้อย่างครบถ้วน
    \item คำนวณหาจำนวนโคโลนีและประเมินค่าชีวมวลของจุลินทรีย์จากตัวอย่างดินเกษตรด้วยวิธีมาตรฐานได้อย่างถูกต้อง
\end{enumerate}

\noindent\textbf{กิจกรรมการเรียนการสอน}
\begin{enumerate}
    \item การบรรยายประกอบสื่อภาพนิ่งและภาพเคลื่อนไหวโครงสร้างเม็ดดิน 3 มิติ
    \item ปฏิบัติการเก็บตัวอย่างดินและตรวจนับจุลินทรีย์ด้วยวิธี Spread Plate บนอาหารจำเพาะ
    \item การศึกษาและอภิปรายผลการตรวจวัดค่า pH ดินต่อปริมาณประชากรเชื้อราและแบคทีเรีย
\end{enumerate}

\noindent\textbf{สื่อการเรียนการสอน}
\begin{enumerate}
    \item เอกสารคำสอนบทที่ 1 และภาพประกอบเวกเตอร์โครงสร้างเม็ดดิน
    \item ตัวอย่างดินจริงจากสวนไม้ผลและแปลงพืชไร่ในจังหวัดจันทบุรี
    \item อาหารเลี้ยงเชื้อ Nutrient Agar (NA) และ Potato Dextrose Agar (PDA)
\end{enumerate}

\noindent\textbf{การประเมินผล}
\begin{enumerate}
    \item การประเมินผลการตรวจนับประชากรจุลินทรีย์ในรายงานผลการทดลองปฏิบัติการ
    \item การทดสอบย่อยประจำบทและการทำแบบฝึกหัดคำนวณชีวมวลท้ายบท
\end{enumerate}

\newpage

% ==========================================================
% เนื้อหาบทเรียน
% ==========================================================

\section{ดินในฐานะระบบนิเวศที่มีชีวิต}
\index{ระบบนิเวศดิน (Soil Ecosystem)}

ดินมิได้เป็นเพียงเศษหิน แร่ธาตุ หรือสารอนินทรีย์ที่ปราศจากชีวิต หากแต่เป็นระบบนิเวศที่มีความซับซ้อนและมีชีวิตชีวามากที่สุดแห่งหนึ่งบนโลก ในดินที่มีความอุดมสมบูรณ์เพียง 1 ช้อนชา (ประมาณ 1 กรัม) สามารถพบประชากรจุลินทรีย์ได้มากกว่าหนึ่งพันล้านถึงหนึ่งหมื่นล้านเซลล์ ซึ่งประกอบด้วยแบคทีเรียนับหมื่นสปีชีส์ สายใยรานับร้อยเมตร และโพรโทซัวอีกนับแสนตัว

สิ่งมีชีวิตขนาดเล็กเหล่านี้ทำหน้าที่เป็นตัวขับเคลื่อนหลักในกระบวนการเปลี่ยนแปลงรูปสาร (Biochemical Transformation) การหมุนเวียนธาตุอาหาร และการสร้างความอุดมสมบูรณ์ของดิน ปริมาณชีวมวลของจุลินทรีย์ในดิน (Soil Microbial Biomass) คิดเป็นสัดส่วนประมาณ 1--5\% ของปริมาณอินทรียวัตถุทั้งหมดในดิน แต่กลับมีกิจกรรมทางชีวเคมีที่ส่งผลกระทบต่อความอุดมสมบูรณ์ของดินและการเจริญเติบโตของพืชอย่างมหาศาล

\section{กลุ่มจุลินทรีย์หลักในดินและบทบาทเชิงนิเวศ}
\index{กลุ่มจุลินทรีย์ในดิน}

ชุมชนจุลินทรีย์ในดินสามารถจัดจำแนกออกเป็นกลุ่มทางชีววิทยาที่สำคัญได้ดังนี้

\subsection{แบคทีเรียในดิน (Soil Bacteria)}
\index{แบคทีเรียในดิน (Soil Bacteria)}
แบคทีเรียเป็นกลุ่มจุลินทรีย์ที่มีจำนวนประชากรมากที่สุดในดิน โดยทั่วไปพบได้ตั้งแต่ $10^7$ ถึง $10^9$ เซลล์ต่อกรัมดินแห้ง มีความหลากหลายทางเมแทบอลิซึมสูงมาก สามารถแบ่งตามโครงสร้างผนังเซลล์และการตอบสนองต่อสีย้อมแกรมได้แก่
\begin{itemize}
    \item \textbf{แบคทีเรียแกรมบวก (Gram-Positive Bacteria)} มักมีผนังเซลล์เปปทิโดไกลแคน (Peptidoglycan) ที่หนา มีความทนทานต่อสภาวะแห้งแล้งและอุณหภูมิสูงได้ดี กลุ่มที่สำคัญอย่างยิ่งในการเกษตรคือจีนัส \textit{Bacillus} ซึ่งสามารถสร้างเอนโดสปอร์ (Endospore) พักตัวได้เป็นเวลานาน และผลิตเอนไซม์ย่อยสลายสารอินทรีย์หลายชนิด
    \item \textbf{แบคทีเรียแกรมลบ (Gram-Negative Bacteria)} มีเยื่อหุ้มชั้นนอก (Outer Membrane) ห่อหุ้มผนังเซลล์บาง มีการตอบสนองต่อสารหลั่งจากรากพืชได้อย่างรวดเร็ว เช่น จีนัส \textit{Pseudomonas}, \textit{Rhizobium}, \textit{Azotobacter}, และ \textit{Burkholderia} ซึ่งหลายชนิดจัดเป็นแบคทีเรียส่งเสริมการเจริญเติบโตของพืช (PGPR)
\end{itemize}

\subsection{แอกทิโนมัยซีท (Actinomycetes)}
\index{แอกทิโนมัยซีท (Actinomycetes)}
แอกทิโนมัยซีทเป็นแบคทีเรียแกรมบวกสายพันธุ์พิเศษที่เจริญเติบโตเป็นสายใยคล้ายเชื้อรา (Filamentous Bacteria) และสืบพันธุ์โดยการสร้างสปอร์ เป็นกลุ่มจุลินทรีย์ที่สร้างสารระเหยชื่อว่า \textbf{จีออสมิน (Geosmin)} ซึ่งเป็นสารที่ก่อให้เกิดกลิ่นไอหอมของไอดินหลังฝนตก (Petrichor) แอกทิโนมัยซีทมีบทบาทสำคัญในการย่อยสลายสารอินทรีย์โครงสร้างซับซ้อน เช่น ไคติน (Chitin) เซลลูโลส (Cellulose) และลิกนิน (Lignin) สกุลที่พบมากที่สุดคือ \textit{Streptomyces} ซึ่งผลิตสารทุติยภูมิและสารปฏิชีวนะมากกว่า 70\% ของยาปฏิชีวนะที่ใช้ในวงการแพทย์และการเกษตร

\subsection{ฟังไจหรือเชื้อราในดิน (Soil Fungi)}
\index{เชื้อราในดิน (Soil Fungi)}
เชื้อรามีชีวมวล (Biomass) คิดเป็นสัดส่วนที่ใหญ่ที่สุดในดิน เนื่องจากมีโครงสร้างเส้นใยรา (Hyphae) ที่ยาวและแผ่กระจายเชื่อมต่อกันเป็นโครงข่ายไมซีเลียม (Mycelium) เชื้อรามีความสามารถในการทนทานต่อสภาพดินกรด (Acidic Soils) ได้ดีกว่าแบคทีเรีย หน้าที่เด่นของเชื้อราคือการผลิตเอนไซม์ภายนอกเซลล์ (Extracellular Enzymes) เพื่อย่อยสลายลิกโนเซลลูโลสในซากพืช และกลุ่มที่อยู่ร่วมกับรากพืชแบบพึ่งพาอาศัยคือ เชื้อราไมคอร์ไรซา (Mycorrhizal Fungi) ซึ่งช่วยเพิ่มพื้นที่ผิวในการดูดซับน้ำและฟอสฟอรัสให้แก่พืช

\begin{definitionbox}{กลูมาลิน (Glomalin)}
กลูมาลิน เป็นไกลโคโปรตีนที่สร้างขึ้นจากผนังเซลล์ของเชื้อราอาร์บัสคูลาร์ไมคอร์ไรซา (AMF) มีคุณสมบัติเหนียวและไม่ละลายน้ำ ทำหน้าที่เสมือนกาวชีวภาพที่ยึดอนุภาคดินเหนียว ดินทราย และอินทรียวัตถุเข้าด้วยกันเป็นเม็ดดินที่มั่นคง ช่วยปกป้องคาร์บอนในดินไม่ให้ถูกย่อยสลายเร็วเกินไป
\end{definitionbox}

\subsection{อาร์เคียและโพรโทซัว (Archaea and Protozoa)}
\index{อาร์เคีย (Archaea)}
ในอดีตเชื่อว่าอาร์เคียอาศัยอยู่เฉพาะในสภาวะแวดล้อมสุดขั้ว แต่การศึกษาทางชีววิทยาโมเลกุลในปัจจุบันพบว่า อาร์เคียในดินมีบทบาทสำคัญอย่างยิ่ง โดยเฉพาะกลุ่ม \textbf{Thaumarchaeota} ซึ่งทำหน้าที่ออกซิไดซ์แอมโมเนีย (Ammonia-Oxidizing Archaea: AOA) ในดินที่มีปริมาณสารอาหารต่ำและดินที่มีความเป็นกรดสูง ส่วนโพรโทซัว (Protozoa) ทำหน้าที่เป็นผู้บริโภคอันดับแรก (Bacterivores) ที่คอยล่ากินเซลล์แบคทีเรีย ช่วยปลดปล่อยธาตุไนโตรเจนที่ถูกกักเก็บในเซลล์แบคทีเรียให้ออกมาในรูปของไอออนแอมโมเนียม ($NH_4^+$) ที่พืชสามารถดูดซึมไปใช้ประโยชน์ได้ทันที

\section{สถาปัตยกรรมเม็ดดินและแหล่งอาศัยระดับไมโคร}
\index{เม็ดดิน (Soil Aggregates)}

โครงสร้างของดินประกอบด้วยการรวมตัวกันของอนุภาคเดี่ยว (ทราย ทรายแป้ง และดินเหนียว) กลายเป็น \textbf{เม็ดดิน (Soil Aggregates)} ซึ่งสามารถจำแนกออกเป็น 2 ขนาดหลัก ดังนี้
\begin{enumerate}
    \item \textbf{เม็ดดินขนาดใหญ่ (Macroaggregates)} มีขนาดเส้นผ่านศูนย์กลางมากกว่า 250 ไมโครเมตร ($>250\ \mu\text{m}$) เชื่อมยึดกันด้วยสายใยรา รากพืช และเมือกพอลิแซ็กคาไรด์
    \item \textbf{เม็ดดินขนาดเล็ก (Microaggregates)} มีขนาดเส้นผ่านศูนย์กลางน้อยกว่า 250 ไมโครเมตร ($<250\ \mu\text{m}$) ยึดเหนี่ยวกันอย่างแน่นหนาด้วยสารประกอบฮิวมัส สารกลูมาลิน และประจุไฟฟ้าของแร่ดินเหนียว
\end{enumerate}

\begin{figure}[H]
\centering
\begin{tikzpicture}[scale=1.0]
    % Outer boundary of Macroaggregate
    \draw[line width=2pt, color=agriEarth, fill=agriEarth!15, rounded corners=25pt] 
        (0,0) .. controls (1,4) and (5,4.5) .. (7,3.5)
              .. controls (9,2.5) and (8.5,-1) .. (7,-2)
              .. controls (5,-3) and (2,-2.5) .. (0,0);
              
    \node at (4.0, 3.8) [font=\footnotesize\bfseries\color{agriEarth}] {เม็ดดินขนาดใหญ่ (Macroaggregate $> 250\ \mu\text{m}$)};

    % Sand grain (quartz)
    \draw[fill=yellow!40!white, draw=agriEarth!80, line width=1.2pt] (1.8, 1.2) circle (0.9cm);
    \node at (1.8, 1.2) [font=\tiny\bfseries\color{agriEarth}, align=center] {เม็ดทราย\\(Quartz)};

    % Silt and clay clusters
    \draw[fill=orange!30!white, draw=agriEarth!80, line width=1pt] (5.2, 1.8) circle (0.7cm);
    \node at (5.2, 1.8) [font=\tiny\bfseries\color{agriEarth}] {ทรายแป้ง};

    % Microaggregate inside
    \draw[line width=1.5pt, color=agriGreen!80!black, fill=agriMint!40, rounded corners=12pt] 
        (3.0,-0.8) rectangle (5.5, 0.6);
    \node at (4.25, -0.1) [font=\scriptsize\bfseries\color{agriGreen!80!black}, align=center] {เม็ดดินขนาดเล็ก\\(Microaggregate)};

    % Plant root fragment / Organic matter
    \draw[line width=4pt, color=agriEarth!90!black, line cap=round] (2.0, 2.5) .. controls (3.5, 2.0) .. (5.0, 2.8);
    \node at (3.5, 3.0) [font=\tiny\color{agriEarth!90!black}] {เศษซากพืช/ฮิวมัส};

    % Fungal hyphae binding particles
    \draw[line width=1.5pt, color=agriLeaf, dashed] (0.8, 0.2) .. controls (2.5, 0.8) and (3.8, -1.5) .. (6.5, -0.5);
    \draw[line width=1.2pt, color=agriLeaf, dashed] (2.2, 1.5) .. controls (4.0, 1.0) .. (6.8, 1.2);
    \node at (6.2, -0.8) [font=\tiny\bfseries\color{agriLeaf}] {สายใยรา (Hyphae)};

    % Bacterial colonies in pores
    \foreach \coord in {(1.2, -1.0), (1.4, -0.8), (1.1, -0.7), (4.8, -1.6), (5.0, -1.4), (5.2, -1.7)} {
        \fill[agriGold] \coord circle (2.5pt);
    }
    \node at (1.3, -1.4) [font=\tiny\bfseries\color{agriGold!80!black}] {โคโลนีแบคทีเรีย};

    % Water film (meniscus)
    \draw[color=agriBlue, line width=2pt, opacity=0.7] (2.7, 1.2) arc (0:80:0.7cm);
    \node at (3.1, 1.6) [font=\tiny\color{agriBlue}] {ฟิล์มน้ำ};

    % Air pore
    \node at (3.5, -1.8) [font=\tiny\color{slateGrey}] {ช่องอากาศ (Air Pore)};
\end{tikzpicture}
\caption{ภาคตัดขวางจำลองโครงสร้างเม็ดดิน (Soil Aggregate Architecture) แสดงแหล่งอาศัยระดับไมโครของจุลินทรีย์ดิน}
\label{fig:soil_aggregate}
\end{figure}

\section{ปัจจัยที่ควบคุมความหลากหลายของจุลินทรีย์ในดิน}
\index{ปัจจัยสิ่งแวดล้อมในดิน}

การกระจายตัว ปริมาณ และกิจกรรมของจุลินทรีย์ในดินถูกกำหนดโดยปัจจัยทางกายภาพและเคมีหลัก 4 ประการ
\begin{enumerate}
    \item \textbf{ความเป็นกรด-ด่างของดิน (Soil pH)} มีอิทธิพลต่อโครงสร้างประชากรจุลินทรีย์อย่างเด่นชัด โดยแบคทีเรียส่วนใหญ่เจริญได้ดีที่สุดในช่วง pH เป็นกลางถึงด่างอ่อน (6.5--8.0) ในขณะที่เชื้อราสามารถทนทานต่อสภาพดินกรดได้ดี (pH 4.5--6.0) ในดินที่มีค่า pH ต่ำกว่า 5.0 กิจกรรมของแบคทีเรียตรึงไนโตรเจนและแบคทีเรียไนตริฟิเคชันจะลดลงอย่างมีนัยสำคัญ
    \item \textbf{ความชื้นและการถ่ายเทอากาศ (Moisture and Aeration)} น้ำในดินทำหน้าที่เป็นตัวกลางในการละลายและแพร่กระจายสารอาหาร เมื่อดินมีความชื้นที่ระดับความจุความชื้นสนาม (Field Capacity) จุลินทรีย์กลุ่มใช้ออกซิเจน (Aerobes) จะทำงานได้สูงสุด หากเกิดน้ำท่วมขัง ช่องอากาศจะถูกแทนที่ด้วยน้ำ ก่อให้เกิดสภาวะขาดออกซิเจน ซึ่งส่งเสริมกิจกรรมของจุลินทรีย์กลุ่มไม่ใช้ออกซิเจน (Anaerobes)
    \item \textbf{อุณหภูมิ (Temperature)} จุลินทรีย์ดินส่วนใหญ่จัดอยู่ในกลุ่มเมโซไฟล์ (Mesophiles) ซึ่งมีอุณหภูมิที่เหมาะสมต่อการเจริญเติบโตระหว่าง 25--35 องศาเซลเซียส
    \item \textbf{ปริมาณและอัตราส่วนคาร์บอนต่อไนโตรเจน (C:N Ratio)} การใส่วัสดุอินทรีย์ที่มีค่า C:N สูง เช่น ฟางข้าว (C:N ประมาณ 80:1) จะทำให้เกิดภาวะจุลินทรีย์แย่งใช้ไนโตรเจนจากดิน (Nitrogen Immobilization) ส่งผลให้พืชแสดงอาการขาดไนโตรเจนชั่วคราว
\end{enumerate}

\section{กรณีศึกษาวิจัย ความหลากหลายและคุณสมบัติทางชีวภาพของดินสวนผลไม้ในจังหวัดจันทบุรี}
\index{จุลินทรีย์ดินสวนผลไม้จันทบุรี}

ในพื้นที่ภาคตะวันออกของประเทศไทย โดยเฉพาะจังหวัดจันทบุรีซึ่งเป็นแหล่งเพาะปลูกไม้ผลเศรษฐกิจที่สำคัญ เช่น ทุเรียน มังคุด และเงาะ สภาพแวดล้อมในดินมีความจำเพาะจากการจัดการสารเคมีและปุ๋ยเคมีต่อเนื่องยาวนาน การศึกษาของ \textbf{จิรภัทร จันทมาลี และคณะ (Chanthamalee et al., 2016, 2017)} เรื่องความหลากหลายของจุลินทรีย์ที่มีประโยชน์และคุณสมบัติทางชีวภาพของดินสวนผลไม้ในพื้นที่ตำบลพลับพลา อำเภอเมือง จังหวัดจันทบุรี ได้ทำการสำรวจประชากรแบคทีเรีย เชื้อรา และแอกทิโนมัยซีท ควบคู่กับค่าความเป็นกรด-ด่างและอินทรียวัตถุในดิน

ผลการวิจัยชี้ให้เห็นว่า ดินสวนผลไม้ที่มีการจัดการแบบผสมผสานร่วมกับอินทรียวัตถุมีปริมาณประชากรแบคทีเรียและแอกทิโนมัยซีทสูงกว่าสวนที่ใช้สารเคมีเข้มข้นอย่างมีนัยสำคัญ ($p < 0.05$) โดยแบคทีเรียกลุ่มสร้างสปอร์สกุล \textit{Bacillus} และแอกทิโนมัยซีทสกุล \textit{Streptomyces} เป็นกลุ่มเด่นที่มีบทบาทสร้างเอนไซม์ย่อยสลายสารอินทรีย์และผลิตสารยับยั้งเชื้อราก่อโรคในระบบนิเวศดินสวนผลไม้ ข้อมูลเชิงประจักษ์นี้ยืนยันว่าการรักษาสมดุลทางชีวภาพของดินเป็นรากฐานสำคัญต่อเสถียรภาพผลผลิตของไม้ผลเมืองร้อนอย่างยั่งยืน

\begin{worked}{การคำนวณประชากรจุลินทรีย์ในดินด้วยวิธี Dilution Plating}
\textbf{โจทย์} ชั่งตัวอย่างดินเกษตรความชื้น 15\% หนัก 10.0 กรัม นำไปเจือจางในน้ำเกลือปลอดเชื้อปริมาตร 90 มิลลิลิตร (ได้สารละลายเจือจางระดับ $10^{-1}$) จากนั้นทำการเจือจางแบบต่อเนื่องจนถึงระดับ $10^{-6}$ แล้วดูดสารละลายปริมาตร 0.1 มิลลิลิตร มาเกลี่ยลงบนจานอาหาร NA บ่มที่อุณหภูมิ 30 องศาเซลเซียส เป็นเวลา 48 ชั่วโมง พบโคโลนีแบคทีเรียเฉลี่ย 142 โคโลนี จงคำนวณหาจำนวนแบคทีเรียต่อกรัมดินแห้ง (CFU/g dry soil)

\textbf{วิธีทำ}
1. คำนวณน้ำหนักแห้งของตัวอย่างดิน
\begin{equation}
\text{น้ำหนักดินแห้ง} = 10.0 \times \left(1 - \frac{15}{100}\right) = 8.5 \text{ กรัม}
\end{equation}
2. คำนวณจำนวนโคโลนีต่อมิลลิลิตรในระดับความเจือจาง $10^{-6}$
\begin{equation}
\text{CFU/mL} = \frac{\text{จำนวนโคโลนี}}{\text{ปริมาตรที่เกลี่ย (mL)}} = \frac{142}{0.1} = 1.42 \times 10^3 \text{ CFU/mL}
\end{equation}
3. คำนวณจำนวนประชากรต่อกรัมดินแห้งทั้งหมด
\begin{equation}
\text{Total CFU/g dry soil} = \frac{1.42 \times 10^3 \times 10^6 \times 100 \text{ mL}}{8.5 \text{ g}} = 1.67 \times 10^{10} \text{ CFU/g}
\end{equation}
นั่นคือ ตัวอย่างดินเกษตรนี้มีประชากรแบคทีเรียเท่ากับ $1.67 \times 10^{10}$ ซีเอฟยูต่อกรัมดินแห้ง
\end{worked}

\section*{คำถามทบทวนท้ายบทที่ 1}
\addcontentsline{toc}{section}{คำถามทบทวนท้ายบทที่ 1}
\begin{enumerate}
    \item จงเปรียบเทียบโครงสร้างผนังเซลล์และบทบาทเชิงนิเวศระหว่างแบคทีเรียแกรมบวกและแกรมลบในดิน
    \item สารกลูมาลิน (Glomalin) มีบทบาทอย่างไรในการสร้างความเสถียรของเม็ดดินและการกักเก็บคาร์บอน
    \item หากดินมีค่าความเป็นกรดรุนแรง (pH 4.2) จะส่งผลกระทบต่อสัดส่วนประชากรระหว่างเชื้อราและแบคทีเรียอย่างไร
    \item จงอธิบายปรากฏการณ์ Nitrogen Immobilization เมื่อเกษตรกรไถกลบตอซังข้าวที่มีค่า C:N ratio สูง
\end{enumerate}

\clearpage